\chapter{Intrinsic Split-Zero support and the integral first-jet defect
of Collatz}
{\small\sffamily Source: \nolinkurl{repository/collatz_reconstruction/research_program/intrinsic_zero_firstjet_20260915/note.md}.\par}


15 September 2026. Source checkpoint:
\nolinkurl{KokunoYumeto/collatz-workbench}, PR \#2,
\nolinkurl{2a77e3bc9435f181bc4c9c3957224a7d93b4e887}. This is an
additive continuation. It corrects a terminology overlap, not the
previously proved equations: the user's informal account of retained
information at a supported zero is not a definition of an exponent-word
history in the weighted-automata companion.

\section{1. The actual foundation and what is retained}

The foundational scalar object is the supplied split-zero construction

\SourceMath{G(R)=R\sqcup\{\tau\},\qquad e=\operatorname{ofR}(0_R)\ne\tau.}{}

The supported copy has its original ring addition and multiplication;
tau is the semiring additive identity and multiplicative absorber. The
archive's original \nolinkurl{SplitZero.lean} states exactly these
operations. Its path-specific commit history reaches the archived July
4, 2026 commit; the June 26 date in its directory name is not being
asserted as the invention date. The current formalization explicitly
cites that archived file and the chapter-14 linear-join-diagram theorem.
This is the documented source chain recovered here, not a claim to have
recovered all earlier private development.

The exact scalar comparison is

\SourceMath{G(R)\longrightarrow R\times\mathbb B,
\quad \tau\mapsto(0,0),\quad \operatorname{ofR}(r)\mapsto(r,1),}{}

with image \(\{(0,0)\}\cup(R\times\{1\})\) and inverse given by those
same cases. Addition in the second coordinate is OR, multiplication is
AND. Case substitution proves both operations and inverse identities.
Reflection to R alone has zero fibre \(\{\tau,e\}\). The separate
support coordinate prevents that collapse. Every map to an additively
cancellative ring sends e to zero, since e+e=e forces h(e)+h(e)=h(e),
and therefore h(e)=0.

For a G(R)-semimodule M, intrinsic support is \(s(m)=em\). The source
theorem identifies the support semilattice with the additive idempotents
and the fibre over l with the R-module

\SourceMath{M_l=\{m:em=l\},}{}

whose zero is l. For l\textless=k its original transition is
\(m\mapsto m+k\). Conversely a coherent diagram \(V_l,\rho_{lk}\)
reconstructs the disjoint union with

\SourceMath{(l,x)+(k,y)=(l\vee k,\rho_{l,l\vee k}x+\rho_{k,l\vee k}y),
\quad e(l,x)=(l,0),\quad\tau(l,x)=(\bot,0).}{}

These are the source's actual quasi-inverse constructions; they do not
require trajectory words. The general theorem permits a nonzero bottom
fibre. The application below has zero bottom fibre, stated explicitly.

There is an exact qualification to the phrase ``the zero remembers.'' On
any nonzero fibre V\_l the map x -\textgreater{} (l,0) has the whole
V\_l as its fibre. In particular it sends both 0 and a nonzero x to the
same element. That element cannot select a unique prior vector. Keeping
the fibre, its relation submodule and a specified comparison retains the
space of possible representatives and the operation, not an
automatically encoded chronological log. Conversely external tau also
has a computable inverse image once its original map is retained; its
output alone supplies no non-bottom support. These statements follow
from the displayed maps, rather than an assertion that there can be no
relation to path history.

An infinite-dimensional fibre or a ring of formal series can be used in
this construction, but its scalar/module operations and comparison must
be specified. The carrier R disjoint union \{tau\} itself adds no formal
summation rule or numerical infinity to R.

For a two-term cochain map f:C-\textgreater D, the source's killed-class
module is

\SourceMath{\{z\in C^1:f^1z\in d_DD^0\}/d_CC^0
\ \cong\ \ker H^1(f).}{1}

The forward map takes a representative to its original source class; its
inverse takes any representative of that source class. Well-definedness
and inverse identities are exactly the original boundary relations. This
is the retained structure used below.

The early Rees continuation computes a comparison defect and retains a
first variation when specialization makes a boundary-character
difference zero. Its particular trace identity and arithmetic estimates
are not transported to Collatz by vocabulary. The common operation here
is explicit: retain a first-order coefficient and the connecting map of
its exact coefficient sequence. The Collatz calculation below is proved
from its own differential.

\section{2. A support diagram of equations, not a sequence of trajectory
words}

Retain the original odd-return map

\SourceMath{X=\{1,3,5,\ldots\},\qquad
T(n)=\frac{3n+1}{2^{a(n)}},\qquad a(n)=\nu_2(3n+1).}{}

For any declared subset E of X, interpret n in E as the actual original
equation/edge \(e_n:n\to T(n)\). Put

\SourceMath{V(E)=E\cup T(E),\quad C_E^0=\mathbb Z^{(E)},\quad C_E^1=\mathbb Z^{(V(E))}.}{}

Parentheses mean finite sums, even when E is infinite. Define

\SourceMath{J_Ee_n=V_n,\qquad P_Ee_n=V_{T(n)},\qquad d_E=J_E-P_E.}{2}

This is an absolute complex: the actual fixed vertex and loop at 1 are
NOT removed. A declared edge with zero differential, notably e\_1,
remains a source basis vector. In finite calculations every target
outside E remains a separate frontier vertex.

The support semilattice is \(\mathcal P(X)\), with union as join and
empty set as bottom. For E subset F all maps are literal inclusions on
original labelled edges and vertices. They commute with J,P,d.~Its
finite-subset subdiagram provides executable approximations; its full
label X defines the infinite object. Reconstruction by the preceding
source formula gives G(Z)-semimodules and internal supported-zero
differentials. The bottom fibres are zero. The order of entering the
equations into a file has no role in E or these maps.

The earlier seed-carrier diagram assigned to E its forward carrier
\(\Omega_E=\bigcup_{j\ge0}T^j(E)\), including 1 where used by that
construction. The natural comparison from this equation diagram to that
carrier diagram is inclusion on each displayed original edge and vertex.
The carrier contains all those labels and their targets; hence this is a
cochain map and commutes with support inclusion. It is not an assertion
that the finite set of known equations already equals its whole forward
carrier.

\subsection{Exact comparison with the elementary Collatz map}

The odd-return source is the existing source of this workbench, not an
unrecorded change to the conjecture. Every positive integer has unique
form 2\^{}q*n with n odd; q initial divisions reach n.~An odd-return
edge n-\textgreater T(n) expands to the original odd step
n-\textgreater3n+1 followed by exactly a(n) divisions. Grouping at
successive odd vertices is its inverse on those elementary blocks. In
particular the original elementary cycle
1-\textgreater4-\textgreater2-\textgreater1 is the one odd-return loop
at 1.

To compare first jets, assign the factor 1+epsilon to each elementary
odd step and factor 1 to each division. Let h(V\_y) be the finite
division chain from y to its odd part, with zero chain for odd y. Send
an odd-return edge to its odd elementary edge plus (1+epsilon) times its
following division chain. Send its vertex to the same odd vertex. The
reverse map sends any elementary vertex to its odd part, every odd
elementary edge to its odd-return edge, and every division edge to zero.
Telescoping proves both are chain maps, their composition on the
odd-return complex is the identity, and on the elementary complex

\SourceMath{d_\epsilon h=I-\iota\pi,\qquad h d_\epsilon=I-\iota\pi}{}

in the respective vertex and edge degrees. For an odd elementary edge
the second equation is the negative of (1+epsilon) times its division
tail; for a division edge it is that original edge. These check both
cases. All columns are finite. Thus the first-jet invariant is preserved
by this specified expansion and retraction. Assigning a new 1+epsilon
factor to EVERY elementary division would be a different coefficient
assignment, counting m+A instead of m; it is not used here.

\section{3. The first-order coefficient map}

Use the ring of dual integers

\SourceMath{\mathbb D=\mathbb Z[\epsilon]/(\epsilon^2).}{}

It has unique coordinates a+epsilon*b and multiplication
\((a,b)(c,d)=(ac,ad+bc)\). Its split extension distinguishes
\(\operatorname{ofR}(\epsilon)\), e, and tau;
\(\operatorname{ofR}(\epsilon)^2=e\), not tau.

Constant reduction \(r:\mathbb D\to\mathbb Z\), a+epsilon*b
-\textgreater{} a, induces G(r). The complete supported-zero fibre is

\SourceMath{G(r)^{-1}(e)=\{\operatorname{ofR}(\epsilon b):b\in\mathbb Z\},
\qquad G(r)^{-1}(\tau)=\{\tau\}.}{3}

Thus already at the scalar level the projected zero has a specified
nontrivial kernel. No trajectory has been stipulated.

On the SAME original graph take

\SourceMath{K_E=[C_E^0\xrightarrow{d_E}C_E^1],\qquad
K_{E,\epsilon}=[C_E^0\otimes\mathbb D\xrightarrow{d_{E,\epsilon}}C_E^1\otimes\mathbb D],}{}
\SourceMath{d_{E,\epsilon}=J_E-(1+\epsilon)P_E=d_E-\epsilon P_E.}{4}

For integral original edge vectors b,c,

\SourceMath{d_{E,\epsilon}(b+\epsilon c)=d_Eb+\epsilon(d_Ec-P_Eb).}{5}

Multiplication by epsilon and constant reduction give the degreewise
exact sequence

\SourceMath{0\longrightarrow K_E\xrightarrow{\epsilon}K_{E,\epsilon}
\xrightarrow r K_E\longrightarrow0.}{6}

Here K\_E is a D-module through r. Equation (5) proves both chain
equations. These maps commute with every support transition.

There is also an exact comparison with the existing weighted companion,
without identifying its term ``history'' with Split-Zero support. On its
ABSOLUTE coefficient ring R=Z{[}u,v{]}, take

\SourceMath{\alpha:R\to\mathbb D,\qquad u\mapsto1+\epsilon,\quad v\mapsto1.}{7}

It sends its edge \(V_n-u v^{a(n)-1}V_{T(n)}\) to (4). Its kernel is
\((v-1,(u-1)^2)\); indeed \(\alpha(f)=f(1,1)+\epsilon\partial_uf(1,1)\),
and successive polynomial remainders in v-1 and u-1 prove both
inclusions. It is surjective, with section of abelian groups a+epsilon*b
-\textgreater{} a+(u-1)b. G(alpha) preserves the two zero labels; its
supported-zero fibre is the supported copy of that ideal. The raw
original endpoints and exponent a(n) remain in the graph record. Formula
(7) is a particular comparison out of the weighted presentation, not the
definition of intrinsic support.

The map (7) does not extend to a unital map from all of
Z{[}v{]}{[}{[}u{]}{]} to D: 1-u is a unit in that source, with its
geometric-series inverse, while its prescribed image -epsilon is not a
unit. The same obstruction holds after G because a supported inverse
would have to multiply -epsilon to 1. The successful replacement used
here is the finite first-jet ring (7) and the original finite-support
complex, not evaluation of that completed inverse.

\section{4. The actual connecting map and its entire zero fibre}

The connecting map of (6) is

\SourceMath{\beta_E:H^0(K_E)\to H^1(K_E),\qquad
c\longmapsto[-P_Ec],\quad d_Ec=0.}{8}

Indeed lift c as a constant vector in K\_epsilon. Its differential is
-epsilon*Pc; removing the epsilon factor gives the displayed class. All
signs are fixed by (4).

Define the integral first-jet defect to be the ACTUAL comparison kernel

\SourceMath{\mathfrak B_E=\ker\bigl(H^1(r):H^1(K_{E,\epsilon})\to H^1(K_E)\bigr).}{9}

\subsection{Theorem 1. Explicit presentation of the retained first-jet
zero fibre}

There are natural isomorphisms

\SourceMath{\boxed{
\mathfrak B_E\cong\operatorname{coker}\beta_E
\cong C_E^1/(d_EC_E^0+P_E\ker d_E)
\cong\epsilon H^1(K_{E,\epsilon}).
}}{10}

\textbf{Proof with inverse maps.} Send the class of y in the middle
quotient to {[}epsilon*y{]}. An original relation d\_Ec maps to
d\_epsilon(epsilon*c). For a cycle b, epsilon*P\_Eb=-d\_epsilon b. Hence
the map is well-defined.

Every class in (9) has representative z0+epsilon*z1 with z0=d\_Eb.
Subtracting d\_epsilon b leaves epsilon*(z1+P\_Eb). Send the class to
{[}z1+P\_Eb{]}. Another choice of b differs by an original cycle, giving
exactly a P\_E*ker d\_E relation. Changing the representative by
d\_epsilon(a+epsilon*c) replaces b by b+a and changes z1+P\_Eb by d\_Ec.
Thus this inverse is well-defined. Both compositions are identity by the
subtraction just performed. Finally every epsilon multiple is
represented by epsilon*y, giving the last equality. QED.

In the Split-Zero source's notation the killed representatives and old
relations are exactly

\SourceMath{K_{\rm rep,E}=d_EC_E^0+\epsilon C_E^1,
\quad B_{\rm old,E}=d_{E,\epsilon}(C_E^0\otimes\mathbb D).}{11}

Their quotient is (9); equation (10) computes it. The supported map is
\((E,[z])\mapsto(E,0)\) on this kernel. At a nonempty label that is not
the external absent point. The maps (10) commute with the original
equation inclusions, so reconstruct a split-linear equivalence of the
full diagrams. The additive and scalar identities are checked fibrewise
after transport to E union F, not by treating the total as an abelian
group with a single zero.

\section{5. Component calculation, including infinite source sets}

For the directed graph with declared edges E, every vertex has at most
one outgoing edge. The free abelian degree-one quotient has one
generator eta\_A for each weak component A. The map is the component
augmentation

\SourceMath{[\sum a_nV_n]\longmapsto(\sum_{n\in A}a_n)_A.}{12}

For its inverse choose the least original integer s\_A in each nonempty
component and map eta\_A to {[}V\_sA{]}. A finite undirected path from n
to s\_A, with a plus or minus sign for each directed original edge, has
boundary V\_n-V\_sA. This proves both compositions and that its kernel
is exactly the ORIGINAL edge-relation image. Finite sums meet only
finitely many components.

The kernel of d\_E has one free integral generator

\SourceMath{c_C=\sum_{n\text{ on }C}e_n}{13}

per actual primitive directed cycle C in E. To prove spanning, restrict
any kernel vector to its finite edge support. A tree leaf forces the
coefficient on its incident edge to vanish. Removing such edges leaves
cycles. Outdegree at most one makes every remaining undirected cycle
directed, and its boundary equation forces all its coefficients to be
equal. Independence follows from disjoint cycle edges. Each weak
component has at most one directed cycle: any two adjacent vertices have
forward paths that meet (or the same finite sink); this property is
preserved across a finite connecting path, and two different cycles
cannot have meeting forward paths.

A component with a cycle of m distinct vertices has

\SourceMath{\beta_E(c_C)=-m\eta_C.}{14}

Every component without a cycle receives no generator under beta. This
proves beta is injective over Z, including for infinite E, and proves
the full defect classification

\SourceMath{\boxed{
\mathfrak B_E\cong
\bigoplus_{C\text{ cycles in }E}\mathbb Z/m_C\mathbb Z
\ \oplus\
\bigoplus_{A\text{ components without a cycle}}\mathbb Z.
}}{15}

Every coordinate map is (12), with reduction modulo m in a cyclic
component; the inverse sends the coordinate generator to
{[}epsilon*V\_sA{]}. Its exact order is m by (12)-(14). A period-one
component contributes a zero GROUP with its original nonempty module
diagram and all equations retained. It is not relabelled absent.

For finite E, a component without a cycle has one frontier vertex with
no declared outgoing equation. Its Z summand is unresolved under that
finite set of equations, NOT evidence of an infinite Collatz orbit. For
E=X, all outgoing equations are present. A component without a cycle has
an injective forward ray and supplies the infinite-component summand in
(15).

\section{6. The complete first-jet cohomology modules}

The preceding defect is not merely a count. On a cycle component of
primitive length m,

\SourceMath{\boxed{H^1(K_{\epsilon}|_C)\cong\mathbb D/(m\epsilon).}}{16}

Choose an original cycle point c.~For every vertex n in its component
let r(n) be the first hitting length to c, with r(c)=0. The map sends
V\_n to (1+epsilon)\^{}\{r(n)\} in the quotient. The relation at c is
\nolinkurl{1-(1+epsilon)^m=-m*epsilon;} every other original edge
respects its first hitting length. The inverse sends 1 to {[}V\_c{]}.
The actual finite path to c gives
{[}V\_n{]}=(1+epsilon)\^{}\{r(n)\}{[}V\_c{]}; the weighted finite cycle
gives the single relation m*epsilon*V\_c=0. These identities prove both
inverse compositions. All paths used for individual vertices are finite,
even with infinitely many incoming vertices.

For a full nonperiodic component choose its actual forward ray x\_j from
its least vertex. A vertex n meeting x\_j after r steps maps to
(1+epsilon)\^{}\{r-j\}. Negative powers are explicitly 1+k*epsilon for
integer k, because (1+epsilon)\^{}\{-1\}=1-epsilon. A later meeting adds
the same number to r and j. The inverse sends 1 to {[}V\_x0{]}; original
edges give both inverse identities. Thus

\SourceMath{\boxed{H^1(K_{\epsilon}|_A)\cong\mathbb D,\qquad
\mathfrak B_A=\epsilon\mathbb D\cong\mathbb Z.}}{17}

The same D presentation holds for a finite open component by choosing
its frontier as reference and using the finite path to it. That
statement has a different declared support and does not turn it into a
nonperiodic component of the full graph.

As abelian groups the cycle module in (16) is Z plus (Z/mZ)*epsilon.
Constant reduction forgets exactly the second term. The base loop m=1
therefore has H1=Z and zero first-jet kernel; a longer cycle has the
retained torsion kernel Z/mZ; a full nonperiodic component has the free
kernel Z. The original constant source spaces remain present in all
three cases.

\section{7. The integral first-jet restatement of Collatz}

A positive fixed point obeys

\SourceMath{(2^a-3)n=1.}{}

For a=1 its solution is negative, for a=2 it is n=1, and for
a\textgreater=3 the positive factor 2\^{}a-3 is at least 5, so no
positive integer solves the equation. The actual positive fixed point is
therefore exactly 1.

On the FULL original absolute graph, (15) now gives

\SourceMath{\boxed{
\mathfrak B_X\cong
\bigoplus_{C\text{ nontrivial positive cycles}}\mathbb Z/m_C\mathbb Z
\ \oplus\
\bigoplus_{A\text{ nonperiodic components}}\mathbb Z.
}}{18}

All periods in the first sum are at least 2. The actual fixed loop
contributes Z/1Z=0 through the unit map -1, without deleting its edge,
vertex, or source fibre.

Consequently the following are equivalent, with the proofs provided by
(10)-(18):

\SourceMath{\boxed{
\begin{gathered}
\text{Every positive odd integer reaches }1;\\
\ker H^1(r:K_{X,\epsilon}\to K_X)=0;\\
\epsilon H^1(K_{X,\epsilon})=0;\\
\beta_X:H^0(K_X)\to H^1(K_X)\text{ is an integral isomorphism}.
\end{gathered}}}{19}

Surjectivity of beta suffices because injectivity was proved in (14).
This is a refinement by a specific comparison kernel, not a claim that
its vanishing is established. It asks that the particular infinitesimal
sources epsilon*V\_n already be original finite boundaries. Their images
under constant reduction are always supported zero, whether or not those
boundaries exist. This is precisely why retaining the original kernel
and relation map matters.

No global endpoint follows from the algebraic classification:
determining that the nonbase component indexing sets are empty remains
the arithmetic content. No new numerical cycle bound is asserted here.

\section{8. Finite witnesses are extractable, not just formal zeros}

The exact certificate for epsilon*V\_n to be an original boundary is a
pair of finite INTEGRAL original edge vectors b,c with

\SourceMath{d_Xb=0,\qquad d_Xc-P_Xb=V_n.}{20}

It is equivalent to \(d_{X,\epsilon}(b+\epsilon c)=\epsilon V_n\), by
(5).

\subsection{Theorem 2. Finite integral witness and actual path
extraction}

Every pair (20) yields an actual finite Collatz path from n to 1 using
only its finite edge support. Conversely, an actual first-arrival path p
from n to 1 yields (20) with

\SourceMath{b=-e_1,\qquad c=\sum_{e\text{ in }p}e.}{21}

\textbf{Proof.} For the converse, d\_Xc=V\_n-V\_1 and -P\_Xb=V\_1; also
d\_Xe\_1=0. This includes n=1 with c=0.

For extraction form the finite graph on the support of b and c,
retaining all its actual targets. In the component containing n, b is an
integer multiple k of its unique cycle sum, or zero when it has no
cycle. Component augmentation of (20) says -k*m=1. Thus a cycle exists
and m=1. The fixed-point calculation proves its vertex is 1. In a finite
functional component with a cycle, every vertex has a forward route to
that cycle: a different terminal frontier or cycle would contradict the
meeting-path property proved in Section 5. Following the declared
original edges therefore reaches 1 within the number of vertices of this
finite graph. The checker executes this extraction after verifying both
integral equations; it does not infer it from a tolerance or a scalar
zero. QED.

This theorem also computes what rationalization would discard. On a
period-m component, m*epsilon*V\_c is a boundary via minus the finite
weighted cycle. Dividing that certificate by m over Q would erase its
order-m class. Over Z, (20) rejects such a division for a singleton
unless m=1.

For the explicitly changed map T\_-(n)=(3n-1)/2\^{}\{nu\_2(3n-1)\}, the
actual cycle is 5-\textgreater7-\textgreater5. Its first-jet defect is
Z/2Z and

\SourceMath{d_\epsilon\bigl(-(e_5+e_7)-\epsilon e_7\bigr)=2\epsilon V_5.}{22}

The checker rejects its half as a nonintegral singleton witness. T\_- is
a control, related to the original SIGNED map by T\_-(n)=-T\_+(-n); it
is not a positive 3n+1 counterexample.

On the original map, the edges 3-\textgreater5-\textgreater1 yield

\SourceMath{d_\epsilon(-e_1+\epsilon(e_3+e_5))=\epsilon V_3.}{23}

The relation involving the supported zero of the fixed loop is an actual
part of this certificate, not a deleted row.

\section{9. Support transitions, repeated words, and original
arithmetic}

For E subset F, the map on (15) sends the generator for an old component
to the generator of its containing new component, with coefficient 1,
followed by the target's modulus. It respects every old relation: a
retained actual cycle remains the same cycle, so its period does not
change; two different cycles cannot later join in a functional graph.
Open components may merge or enter a known cycle. The matrices compose
exactly, including modulus-one labels. This is a diagram indexed by
equation support, not by chronological time.

Every finite vector and every relation in the full construction uses
finitely many equations. Thus

\SourceMath{\mathfrak B_X\cong\mathop{\rm colim}_{E\subset X,\ E\text{ finite}}\mathfrak B_E.}{24}

Direct proof: represent a class by epsilon times finitely many original
vertex vectors, include edges from those vertices in E, and use (10).
Any equality in the full quotient has an original finite d-vector plus a
finite cycle vector; putting their supports in a larger F proves the
equality at that finite stage. This gives both inverse descriptions of
the colimit. It agrees with the standard exactness of filtered module
colimits, without assuming compactness of a space of infinite integer
orbits.

A declared position word is not automatically an actual primitive cycle.
An actual solution x\_i of 2\^{}\{a\_i\}x\_(i+1)-3x\_i=1 gives the
position-to-orbit maps e\_i -\textgreater{} e\_(x\_i), V\_i
-\textgreater{} V\_(x\_i). They commute with both J and P, hence with
d\_epsilon and beta. A word traversing an actual primitive cycle r times
maps its position cycle to r times the actual cycle. In particular the
position circle for (2,2) has a Z/2Z first-jet defect, but its original
realization x=(1,1) maps both edges to e\_1, and maps that defect to
Z/1Z=0. The kernel is accounted for by this explicit map. Using position
length in place of actual primitive length would give a false
counterexample.

The previous integral forcing class {[}C{]} in Z/(2\^{}A-3\^{}m)Z
remains a separate retained invariant connected through this actual
realization map. Integrality, positivity and exact valuations are
verified before a position word is assigned an actual cycle. No equality
of this forcing modulus with the first-jet index m is asserted. The
numerical endpoints themselves recover the original exponent by
2\^{}a=(3n+1)/T(n); its labelled edge is not replaced by an unmarked
abstract cycle.

\section{10. Coefficient and infinity comparisons}

From (18), coefficient extension to Q kills exactly the finite cyclic
defect summands and maps each free nonperiodic summand injectively to Q.
Thus rational vector-space dimensions alone would miss all finite
cycle-index obstructions in this particular comparison. Split-Zero
support is preserved by G(Z)-\textgreater G(Q), but the original
integral kernel is still required. This is a computed loss, not a claim
of an absence of relations between the coefficient settings.

For a prime p and integer r\textgreater=1, the exact quotient
presentation gives

\SourceMath{\mathfrak B_X\otimes\mathbb Z/p^r
\cong\bigoplus_C\mathbb Z/\gcd(m_C,p^r)\mathbb Z
\oplus\bigoplus_A\mathbb Z/p^r\mathbb Z.}{25}

Every nontrivial finite m has a prime divisor, so vanishing of these
reductions for all primes is equivalent to vanishing of (18). That
detection claim uses the PROVED direct-sum form (18); it is not claimed
for arbitrary abelian groups.

The infinite graph in (18) is handled without infinite sums of
incident-edge coefficients. All chain columns remain finite. Along a
specified injective original ray
x\_0-\textgreater x\_1-\textgreater\ldots, an enlarged module allowing
locally finite edge chains contains epsilon*sum\_j e\_(x\_j), whose
componentwise boundary is epsilon*V\_(x\_0). Each vertex of that ray has
at most two contributing edges, so this calculation is well-defined. It
displays a free source class killed by that enlarged comparison. It does
not compute the entire kernel of every locally finite completion.
Retaining the original finite-support relation submodule prevents this
completed primitive from being accepted as (20).

\section{11. What the source rereading changes in the research
programme}

The scalar core is not a device that automatically stores words. The
useful research object is a specified support diagram and a specified
map, with its full fibre of representatives carried to supported zero.
The original Rees work demonstrates that a specialization may erase a
zeroth-order observation while a first-order comparison retains a
defect. Sections 3-7 implement that operation directly on the original
Collatz equations.

The newly exposed arithmetic target is the vanishing of the actual
first-jet kernel (9) at full support, or equivalently the integral
surjectivity in (19). It separates, within the SAME comparison, finite
recurrence index from nonperiodic persistence and retains the fixed-loop
source rather than starting by deleting it.

There are two concrete proof-development tasks, neither assumed here:

\begin{itemize}
\tightlist
\item
  Construct complete arithmetic families of finite pairs (20), retaining
  their source parameters and original valuations, and prove their
  coverage of all original positive sources. The small checker can
  verify individual instances and extracted paths; universal family
  coverage requires a written proof, not a sample or timeout. Equation
  (24) identifies the precise finite-support exhaustion that such a
  proof must establish.
\item
  Control the original nonzero defect classes with arithmetic
  obstructions. The torsion part requires excluding actual primitive
  periods m\textgreater=2, not merely generic matrix cycles. The free
  part requires excluding an original nonperiodic component. Prime
  tests, tensor amplification, or geometric estimates are useful only
  with maps returning to this same original integral kernel and a proof
  of their coverage or bound. No estimate from the Riemann-hypothesis
  programme is imported.
\end{itemize}

A finite support's open free generator is a question to extend, not a
divergence verdict. Example: E=\{3,13\} has one open component at 5;
adding the edge at 5 reaches an open vertex 1; adding the ACTUAL
fixed-loop equation makes its first-jet defect zero. The source modules
and all previously entered equations remain in the same diagram
throughout.

This is a new exact presentation and executable witness interface, not a
claimed proof of the global arithmetic vanishing. It gives no new
numerical cycle bound, no new Lean execution, and no independent
external mathematical audit.

\section{12. Reproduction and scope}

\nolinkurl{firstjet.py} constructs the declared equation diagram, both
original matrices, its full first-jet cohomology defect presentation,
support-transition matrices and integral singleton certificates.
\nolinkurl{verify.py} checks 128 finite original equation sets; 2,187
support inclusions; 1,024 support-composition diagrams; abstract
position-circle controls of lengths 1 through 32; the actual -1-forcing
cycle; and 1,024 original positive odd singleton certificates, with
25,437 separately replayed original return equations. Twelve malformed
or false inputs are rejected. All tests remain active under optimized
Python.

The separate integer matrix calculation constructs unimodular row and
column maps U,V and checks U*A*V=D, det U=det V=+/-1, and the actual
diagonal divisibility. Its row additions, column additions, swaps and
sign changes have the explicit inverse elementary operations. It does
not substitute the diagonal matrix for the retained original source.
Prime reductions of the original matrices provide additional independent
rank checks.

The general statements, including full infinite source (18), are proved
above. Finite matrices and path replays are implementation regression
evidence. Ordinary and optimized runs each passed 39,211 exact checks
and produced the same recorded bytes. \nolinkurl{replay.py} verifies the
source manifest and both executions.

\section{13. Source locators and attribution}

\begin{itemize}
\tightlist
\item
  {[}S0{]} Archived scalar core:
  \nolinkurl{KokunoYumeto/modern-latex-manuscripts}, revision
  \nolinkurl{f7ff59b176c7dc3941babd4cb9272dffc653070d},
  \nolinkurl{formalization/lean/classical_candidates_20260626/split_support_sidecar/SplitZero.lean},
  blob \nolinkurl{ff991f7383922e71cdf0e4a3bc85e89e18f808ef}. Read in
  full. Its retrieved path history has the archived July 4 commit
  \nolinkurl{d7def24933004584b687842fe75e92ed5ee655a1}.
\item
  {[}S1{]} Chapter 14,
  \nolinkurl{satellites/14_globalization_cue_split_zero_and_shifted_sheets.tex},
  Zeta revision \nolinkurl{42d00e359b16d52ca71568ce5e3db5341949d929},
  blob \nolinkurl{49318a578b96462384000afa16794cfd5fe492c1}: scalar
  coordinate formulas, intrinsic support and linear-join-diagram
  equivalence. Read the algebraic source through its ideal/zero-divisor
  sections. Its historical presentation and empty-join qualifications
  are corrected in {[}S2{]}; those historical errors are not adopted
  here.
\item
  {[}S2{]} \nolinkurl{formal/splitzero/README.md} and
  \nolinkurl{MATHEMATICAL_NOTE.md}, Zeta revision
  \nolinkurl{652d3c7e269cef11e0f667be0401abb6014314c2}, blobs
  \nolinkurl{4279feaeb3a80f203684d7f6b00a4e208b71d92d} and
  \nolinkurl{7531d36aa44d8213e0e80d15f79668b6c8839ce5}. Read in full,
  including the required {[}tau{]}=0 presentation relation and the
  non-bottom ideal-poset correction.
\item
  {[}S3{]} \nolinkurl{formal/splitzero/DERIVED_MATHEMATICS.md}, same
  Zeta revision, blob
  \nolinkurl{fe05d32ccd314daac7c59d225c9cac3d4c1bf3d3}, Sections 1-6:
  reconstruction, internal quotient, killed-representative kernel, and
  actual homology maps. The earlier source's own Lean scope is not a
  Lean verification of this new note.
\item
  {[}S4{]}
  \nolinkurl{workbenches/split-support-rees-trace/RESEARCH_NOTE.md},
  same Zeta revision, blob
  \nolinkurl{168775d2be296d74b8243cebd3e9fc8942f2eb31}, Sections 1-5.
  The retained first-variation calculation motivated the distinct
  comparison proved here; its tensor-growth estimates are not reassigned
  to Collatz.
\item
  {[}C{]} Preserved Collatz predecessor:
  \nolinkurl{supported_structure_20260915/note.md}, revision
  \nolinkurl{2a77e3bc9435f181bc4c9c3957224a7d93b4e887}, blob
  \nolinkurl{115af68c1684147e840cdd6cf690269490058f95}; the uploaded
  copy matches this blob. Its automata/word use of ``history'' remains
  that named presentation, not an intrinsic definition of Split-Zero
  cohomology.
\item
  Standard homological constructions: Stacks Project, cones and termwise
  split sequences, \nolinkurl{https://stacks.math.columbia.edu/tag/014D}
  ; Tor exact sequences,
  \nolinkurl{https://stacks.math.columbia.edu/tag/00M0} ; exact filtered
  module colimits, \nolinkurl{https://stacks.math.columbia.edu/tag/00DB}
  . Equations (8)-(11) and (24) include direct proofs rather than
  relying on an unspecified connecting-map convention.
\end{itemize}

Research direction and Split-Zero framework: KokunoYumeto. This note's
Collatz first-jet application, proof exposition and executable checks
were developed in this continuation. No global priority claim is made
for the standard scalar, homological or graph-theoretic operations. No
Zeta file, inherited contribution, frozen edition, or remote main branch
is changed by this addition.
